% !TEX root = ../../../main.tex
% Block card: detection.summation_threshold -- "Summation threshold" (Part I, the blocks chapter).
\cardfield{What it does}
Reads the activity curve from \blockref{detection.wavelet_summation} together with the signal
itself, and marks the stretches the method of \citet{dehshibi2021electrical} calls spikes.

\cardfield{When to use it}
As the last block of the \texttt{dehshibi\_spikes} template. It will also take any other block
that gives one score per sample, so you can try a different activity curve in front of the
same spike rules.

\cardfield{What goes in, what comes out}
In goes a score curve (and, automatically, the signal). Out come the spikes, one stretch each,
scored by how far the signal moves inside it. The block's page also shows each step on the way
(a ``funnel''): the candidate stretches, those that passed, those set aside, and the final
spikes.

\cardfield{Settings that matter}
\emph{epsilon\_factor} (0.05) is how prominent a rise or dip in the score must be, as a share
of the score's range within the window: raise it for fewer candidates. \emph{min\_separation\_s}
(60\,s) is the shortest gap between those rises and dips, and the shortest candidate kept.
\emph{window\_s} must match the one used by \blockref{preprocessing.wavelet_transform}.

\cardfield{Watch out for}
Every threshold is relative to its own window, so a quiet window still yields candidates. On
the M2\_aug recording the method put a spike in 91\,\% of the ten-minute windows a researcher
had marked interesting, but also in 86\,\% of those marked not interesting, and half of the
time lay inside a spike. Treat it as a faithful baseline, not a selective detector.
